% 70-69(70쪽) — 도함수 y=f'(x) 의 그래프(구간 [a, b]). 스캔 화질이 낮아 TikZ 로 다시 그림.
% 원문에 있는 것만: 곡선 · 안내 점선 넷(x_1, x_2, x_4, x_5) · 라벨 a, x_1, x_2, x_3, O, x_4, x_5, b, x, y, y=f'(x).
% 원문에 눈금·좌표값은 없다. f'(x)=0 이 되는 자리(a, x_3, b)와 극값 자리만 원문 그대로 옮긴다.
\begin{tikzpicture}[x=0.40cm, y=0.36cm, line width=0.5pt]
  \draw[->, >=stealth, line width=0.7pt] (-7.3,0) -- (5.25,0) node[below=1pt, font=\footnotesize] {$x$};
  \draw[->, >=stealth, line width=0.7pt] (0,-3.2) -- (0,3.85) node[left=1pt, font=\footnotesize] {$y$};
  \draw[densely dashed, line width=0.35pt] (-4.97,0) -- (-4.97,2.14);
  \draw[densely dashed, line width=0.35pt] (-2.24,0) -- (-2.24,2.57);
  \draw[densely dashed, line width=0.35pt] (1.14,0) -- (1.14,-0.57);
  \draw[densely dashed, line width=0.35pt] (2.29,0) -- (2.29,-1.96);
  \draw[line width=0.6pt] plot[smooth, tension=0.6] coordinates {
    (-6.70,-1.70) (-6.45,-0.95) (-6.21,0.00) (-5.95,0.72) (-5.60,1.45) (-5.30,1.90)
    (-4.97,2.14) (-4.65,2.00) (-4.30,1.55) (-4.00,1.15) (-3.71,0.90) (-3.45,0.98)
    (-3.10,1.30) (-2.75,1.85) (-2.45,2.35) (-2.24,2.57) (-2.05,2.45) (-1.85,2.10)
    (-1.60,1.30) (-1.29,0.00) (-1.05,-0.95) (-0.75,-1.75) (-0.45,-2.20) (-0.15,-2.38)
    (0.05,-2.39) (0.40,-2.30) (0.70,-1.95) (0.95,-1.05) (1.14,-0.57) (1.35,-0.68)
    (1.60,-1.05) (1.90,-1.55) (2.10,-1.82) (2.29,-1.96) (2.50,-1.90) (2.75,-1.58)
    (3.05,-0.98) (3.30,-0.44) (3.47,0.00) (3.75,0.85) (4.05,1.90) (4.36,2.90) };
  \node[below=1pt, font=\footnotesize] at (-6.03,0) {$a$};
  \node[below=1pt, font=\footnotesize] at (-4.97,0) {$x_1$};
  \node[below=1pt, font=\footnotesize] at (-2.24,0) {$x_2$};
  \node[above=1pt, font=\footnotesize] at (-0.85,0) {$x_3$};
  \node[below left=-1pt and -2pt, font=\footnotesize] at (0,0) {$\pt{O}$};
  \node[above=1pt, font=\footnotesize] at (1.14,0) {$x_4$};
  \node[above=1pt, font=\footnotesize] at (2.29,0) {$x_5$};
  \node[below=1pt, font=\footnotesize] at (3.62,0) {$b$};
  \node[font=\footnotesize] at (3.30,3.30) {$y=f'(x)$};
\end{tikzpicture}
